\section{La sabiduría es nuestra relación con todas las cosas}

Esta sección culmina en la definición central de la segunda mitad de la entrevista: la sabiduría es nuestra relación con todas las cosas. Esta definición es singular porque no entiende la sabiduría como IQ, cantidad de conocimiento o puntaje de razonamiento, sino como la calidad de las relaciones entre personas, entre las personas y las cosas, las organizaciones, la tecnología y el mundo. Por ello, el sentido de la IA tampoco se limita a sustituir a las personas: consiste en liberarlas de acciones de bajo valor que consumen energía, para que dispongan de más energía con la cual atender relaciones de mayor calidad.

\begin{figure}[H]
\centering
\includegraphics[width=0.92\textwidth]{wisdom-relation-map.png}
\caption{La sabiduría es relación: la sabiduría surge del contacto continuo con las personas, las cosas, las organizaciones, la tecnología y el mundo. Diagrama conceptual propio, elaborado a partir de la conversación de 02:26:47--02:31:10.}
\end{figure}

\begin{knowledgebox}{Lectura de la figura: la sabiduría no es un puntaje mental}
Las personas, las cosas, las organizaciones, la tecnología y el mundo constituyen en conjunto la fuente de la sabiduría. Si una persona solo razona dentro de su cabeza, sin entrar en contacto con personas y hechos reales, su sabiduría se empobrece; si una IA solo circula entre tokens, sin pasar a la acción y a la retroalimentación, también le resulta difícil generar valor real.
\end{knowledgebox}

\subsection{El sentido de la IA: reducir las acciones que consumen energía}

El apartado anterior definió la sabiduría como relación; este apartado se ocupa de cómo la IA cambia la distribución de las acciones dentro de esas relaciones. Li Xiang (李想) pone el ejemplo del personal de ventas: a los buenos vendedores y especialistas de producto les gusta tratar con los clientes, porque ese contacto les da energía; en cambio, las invitaciones repetitivas, los procesos de bajo valor y la comunicación mecánica pueden agotarla. Uno de los valores de la IA es automatizar estas acciones que consumen energía, para que las personas dediquen su energía a tareas con más sabiduría y mayor calidad de relación.

\begin{importantbox}{La relación entre la IA y las personas}
La IA no solo sustituye trabajo: también redistribuye la energía de las personas. Un sistema de IA verdaderamente bueno debería reducir las acciones de bajo valor, repetitivas y desgastantes, y al mismo tiempo conservar o fortalecer el contacto de alta calidad de las personas con los clientes, los productos, las organizaciones y el mundo.
\end{importantbox}

\subsection{Del envoltorio a las capacidades básicas}

Si el sentido de la IA está en asumir más action, al final es inevitable volver a la cuestión de las capacidades básicas. Hacia el cierre, Li Xiang expone su opinión sobre aplicaciones Agent como «Xiaohong» (小红). Considera que llevar la IA hacia la action real, la action profesional y la to-do list es un intento importante; pero si los competidores son empresas con modelos sólidos, no basta con un envoltorio, máquinas virtuales y capacidad de uso de herramientas: las capacidades necesarias del modelo hay que entrenarlas de verdad. Este planteamiento deja una advertencia muy clara para los emprendimientos de aplicaciones de IA: el punto de entrada, los flujos y la experiencia de producto pueden ponerse en marcha primero, pero al final hay que completar las capacidades básicas.

\begin{figure}[H]
\centering
\includegraphics[width=0.94\textwidth]{shell-vs-model-ability.png}
\caption{Envoltorio frente a capacidad del modelo: la innovación en aplicaciones debe avanzar hacia la action y, en última instancia, completar las capacidades básicas. Diagrama conceptual propio, elaborado a partir de la conversación de 02:44:21--02:45:30.}
\end{figure}

\begin{knowledgebox}{Lectura de la figura: el envoltorio tiene valor, pero no hay que quedarse ahí}
A la izquierda, el envoltorio y las herramientas permiten validar con rapidez las necesidades de los usuarios y los flujos de trabajo; a la derecha, las capacidades básicas determinan el techo a largo plazo. Si una aplicación Agent quiere llegar a la action profesional, tiene que completar de manera gradual el modelo, la memoria, la invocación de herramientas, la evaluación y la confiabilidad.
\end{knowledgebox}

\begin{warningbox}{No interpretar «el envoltorio no es innovación» como una postura contra las aplicaciones}
La innovación en aplicaciones es muy importante, porque es la más cercana a los usuarios y a los escenarios de acción. Sin embargo, cuando el escenario pasa a tareas de alto valor, con fuerte exigencia de alineación y alta confiabilidad, depender solo de modelos externos y de flujos de trabajo delgados choca con un techo, y las capacidades básicas tienen que estar a la altura.
\end{warningbox}

\subsection{Resumen de la sección}

«La sabiduría es relación» devuelve todo el episodio de la tecnología a las personas. VLA, Agent OS, los modelos del mundo y la metodología organizacional no buscan, en última instancia, lucirse técnicamente, sino que la inteligencia se integre mejor en la relación entre las personas y el mundo.
